% Options for packages loaded elsewhere
\PassOptionsToPackage{unicode}{hyperref}
\PassOptionsToPackage{hyphens}{url}
%
\documentclass[
  11pt,
  a4paper,
]{article}
\usepackage{amsmath,amssymb}
\usepackage{iftex}
\ifPDFTeX
  \usepackage[T1]{fontenc}
  \usepackage[utf8]{inputenc}
  \usepackage{textcomp} % provide euro and other symbols
\else % if luatex or xetex
  \usepackage{unicode-math} % this also loads fontspec
  \defaultfontfeatures{Scale=MatchLowercase}
  \defaultfontfeatures[\rmfamily]{Ligatures=TeX,Scale=1}
\fi
\usepackage{lmodern}
\ifPDFTeX\else
  % xetex/luatex font selection
  \setmainfont[]{DejaVu Serif}
  \setsansfont[]{DejaVu Sans}
  \setmonofont[]{DejaVu Sans Mono}
\fi
% Use upquote if available, for straight quotes in verbatim environments
\IfFileExists{upquote.sty}{\usepackage{upquote}}{}
\IfFileExists{microtype.sty}{% use microtype if available
  \usepackage[]{microtype}
  \UseMicrotypeSet[protrusion]{basicmath} % disable protrusion for tt fonts
}{}
\makeatletter
\@ifundefined{KOMAClassName}{% if non-KOMA class
  \IfFileExists{parskip.sty}{%
    \usepackage{parskip}
  }{% else
    \setlength{\parindent}{0pt}
    \setlength{\parskip}{6pt plus 2pt minus 1pt}}
}{% if KOMA class
  \KOMAoptions{parskip=half}}
\makeatother
\usepackage{xcolor}
\usepackage[margin=22mm]{geometry}
\usepackage{longtable,booktabs,array}
\usepackage{calc} % for calculating minipage widths
% Correct order of tables after \paragraph or \subparagraph
\usepackage{etoolbox}
\makeatletter
\patchcmd\longtable{\par}{\if@noskipsec\mbox{}\fi\par}{}{}
\makeatother
% Allow footnotes in longtable head/foot
\IfFileExists{footnotehyper.sty}{\usepackage{footnotehyper}}{\usepackage{footnote}}
\makesavenoteenv{longtable}
\setlength{\emergencystretch}{3em} % prevent overfull lines
\providecommand{\tightlist}{%
  \setlength{\itemsep}{0pt}\setlength{\parskip}{0pt}}
\setcounter{secnumdepth}{5}

\usepackage{amsmath,amssymb,booktabs,longtable,array,etoolbox,xurl,fvextra,graphicx}
\graphicspath{{docs/}}
\usepackage{fancyhdr}
\pagestyle{fancy}
\fancyhf{}
\fancyhead[L]{\small Singularity: Relational Heights and Arithmetic Visibility}
\fancyfoot[C]{\thepage}
\setlength{\headheight}{15pt}
\setlength{\emergencystretch}{3em}
\setlength{\tabcolsep}{3pt}
\renewcommand{\arraystretch}{1.12}
\AtBeginEnvironment{longtable}{\footnotesize\everymath{\scriptstyle}}
\fvset{breaklines=true,breakanywhere=true}
\allowdisplaybreaks
\makeatletter
\renewcommand*\l@section{\@dottedtocline{1}{1.5em}{3.6em}}
\renewcommand*\l@subsection{\@dottedtocline{2}{3.8em}{4em}}
\renewcommand*\@pnumwidth{2.5em}
\renewcommand*\@tocrmarg{3.5em}
\makeatother
\AtBeginDocument{\hypersetup{colorlinks=true,linkcolor=teal,urlcolor=teal}}
\ifLuaTeX
  \usepackage{selnolig}  % disable illegal ligatures
\fi
\IfFileExists{bookmark.sty}{\usepackage{bookmark}}{\usepackage{hyperref}}
\IfFileExists{xurl.sty}{\usepackage{xurl}}{} % add URL line breaks if available
\urlstyle{same}
\hypersetup{
  pdftitle={Relational Height Populations and Arithmetic Visibility},
  pdfauthor={Singularity project contributors},
  hidelinks,
  pdfcreator={LaTeX via pandoc}}

\title{Relational Height Populations and Arithmetic Visibility}
\usepackage{etoolbox}
\makeatletter
\providecommand{\subtitle}[1]{% add subtitle to \maketitle
  \apptocmd{\@title}{\par {\large #1 \par}}{}{}
}
\makeatother
\subtitle{Exact examples, nested boundaries, and primewise
singularities}
\author{Singularity project contributors}
\date{18 October 2026}

\begin{document}
\maketitle

\hypertarget{abstract}{%
\section*{Abstract}\label{abstract}}
\addcontentsline{toc}{section}{Abstract}

We study nonnegative integer heights on a finite graph, counted jointly
by total height and quadratic edge energy. Removing a common translation
gives canonical minimum-zero profiles. An additional span mark separates
binary cuts from nested height configurations. We give explicit
seven-vertex examples in which equal Laplacian characteristic
polynomials, or equal complete binary cut-volume censuses, coexist with
unequal height coefficients. The first differences occur at total
heights two and three. For both pairs, the unmarked local coefficient
tables through total height four, equality below the first difference,
the first-difference polynomial and its stated factorization are checked
by the Lean 4 kernel against a semantic definition of the coefficients.
Apart from Lean itself, the only unverified link is a transcription
script with recorded input digests. A third comparison, between the cube
and the eight-vertex Möbius ladder, delays the first difference to total
height four; it is checked by exact software replay only. For a fixed
energy weight q in the closed unit interval, we prove the following. If
the first difference occurs at total height r, it fixes the exponent of
a zeta factor in the ratio of the associated prime products, and hence
the power singularity at s=1/r. For the two seven-vertex pairs at q=1/2,
the regular factor at s=1/r is enclosed in explicit rigorous intervals.
The examples come with standard-library exact software, independent
finite replay, and an exhaustive census of connected graphs through six
vertices. General reconstruction, unbounded delayed visibility and
historical priority are left open.

\hypertarget{question-and-background}{%
\section{Question and background}\label{question-and-background}}

A graph parameter can retain many numerical facts while forgetting the
correspondence between the configurations that produced them. The binary
cut-volume census is a useful example. It counts vertex subsets jointly
by cardinality and boundary size, so equality of this polynomial is
stronger than equality of a single energy moment. Nevertheless, it does
not retain which subsets are nested or which edges cross two of their
boundaries simultaneously. Integer heights provide a direct way to
observe those relationships.

The ingredients have established antecedents. Integral cut lattices
organize graph gradients modulo constants {[}BHN1997{]}. Bivariate Ising
polynomials organize binary configurations {[}AM2009{]}. Integer-valued
Gaussian models on lattices use the same squared nearest-neighbour
energy {[}BPR2024{]}. Formal Euler transforms reorganize coefficient
data into cyclotomic factors {[}Moree2005{]}. We use these ingredients
with explicit marks and a restricted inverse question. The paper does
not claim a new origin for the individual ingredients. Nor does it claim
that the proposed scalar-height observations determine every graph. No
completed priority search is claimed.

There are two logically distinct constructions. First, a local
generating function records a finite graph's integer heights. Second, an
Euler product attaches that same local function to every rational prime.
The first determines graph-information content. The second supplies an
arithmetic organization and a precise comparison of singularities.
Keeping this distinction visible prevents a prime product from being
credited with information that it merely reencodes.

Results labeled SNG-\ldots{} refer to the full proofs in the project
monograph (\texttt{docs/MONOGRAPH.md}). The proofs given here are
complete for the statements as displayed. Lean theorem names refer to
the core Lean 4.12.0 package in \texttt{formal/}.

\hypertarget{local-populations-and-their-marks}{%
\section{Local populations and their
marks}\label{local-populations-and-their-marks}}

Let G be a finite connected simple loopless graph on vertices
\(0,\ldots,k-1\), with \(k\ge2\). Write \(d_i\) for its degrees and
\(A_{ij}\) for its adjacency matrix. Connectedness is used only in
Section 5 and in the certified enclosures of Section 4. Propositions 1
and 2 and Theorem 3 hold without it. For a nonnegative integer vector a,
put

\[
 |a|=\sum_i a_i,\qquad
 E_G(a)=\sum_{\{i,j\}\in E(G)}(a_i-a_j)^2.
\]

The unnormalized series and the fully marked normalized series are

\[
 L_G(x,q)=\sum_{a\ge0}x^{|a|}q^{E_G(a)},\qquad
 \mathcal B_G(x,q,z)=\sum_{\min a=0}x^{|a|}q^{E_G(a)}z^{\max a}.
\]

Every coefficient of x is a finite polynomial in q, because there are
only \(\binom{m+k-1}{k-1}\) nonnegative vectors of total height m
(SNG-F01). Hence these definitions are formal before any analytic
parameter is chosen. Counts are of original labeled assignments; no
quotient by graph automorphisms is taken. For 0 at most q at most one
and absolute value of x below one, comparison with \((1-|x|)^{-k}\) also
gives analytic convergence.

\textbf{Proposition 1 (translation normalization; SNG-F02).} With
\(B_G(x,q)=\mathcal B_G(x,q,1)\),

\[
 L_G(x,q)=\frac{B_G(x,q)}{1-x^k}.
\]

\textbf{Proof.} Every nonnegative vector has the unique form
\(a=b+t\mathbf1\), where \(t=\min a\) and \(\min b=0\). Edge
differences, and hence energy, are unchanged. Total height increases by
kt. Summing the common translation gives the geometric factor. This is a
formal identity with no convergence assumption. \(\square\)

The coefficientwise form of Proposition 1 is kernel-checked for
arbitrary edge lists, including disconnected graphs, as
\texttt{allCount\_eq\_shiftSum} and \texttt{allCount\_rec}
(\texttt{formal/Singularity/Normalization.lean}). The span-marked
identity, with z retained, is not formalized.

The span mark records a genuine distinction. The binary cut polynomial

\[
 I_G(x,q)=\sum_{S\subseteq V(G)}x^{|S|}q^{|\partial S|}
\]

satisfies \(I_G=1+x^k+[z]\mathcal B_G\) (SNG-F11). Thus the full marked
series contains the binary census explicitly. The specialization z=1
merges spans; a general recovery theorem for the binary sector from that
specialization is not established here.

\textbf{Proposition 2 (nested-boundary energy; SNG-F09, SNG-F10).} A
normalized vector of span h is equivalent to a chain of nonempty proper
sets

\[
 S_h\subseteq\cdots\subseteq S_1,\qquad
 a=\sum_{t=1}^h\mathbf1_{S_t}.
\]

Repeated sets are permitted. If \(e(U,W)\) counts edges between disjoint
sets, then

\[
 |a|=\sum_t|S_t|,\qquad
 E_G(a)=\sum_t|\partial S_t|
       +2\sum_{r<t}e(S_t,V\setminus S_r).
\]

\textbf{Proof.} Take \(S_t=\{i:a_i\ge t\}\). Each coordinate is
recovered by counting its occupied thresholds, proving the bijection and
total-height identity. On an edge, every nonzero threshold difference
has the same sign. Expanding the square of their sum gives the
individual threshold crossings and twice every pair of simultaneous
crossings. For \(r<t\), simultaneous crossing means that one endpoint
lies in \(S_t\) and the other outside \(S_r\). Summing edges proves the
formula. \(\square\)

The threshold reconstruction, the volume identity and the energy
identity are kernel-checked for every edge list and every profile as
\texttt{reconstruct}, \texttt{vol\_eq\_levels} and
\texttt{energy\_nested}. Nestedness and properness of the threshold sets
are checked as \texttt{thr\_nested} and
\texttt{thr\_proper\_of\_minZero}, and the span count as
\texttt{span\_eq} (\texttt{formal/Singularity/Thresholds.lean}). The
converse direction of the bijection, from an arbitrary chain to a
profile, is proved above but is not stated in Lean.

For two levels, this becomes
\(E(\mathbf1_S+\mathbf1_T)=|\partial S|+|\partial T|+2e(S,V\setminus T)\),
for nonempty S contained in proper T, allowing equality. The final term
explains why a distribution of isolated cuts can miss the multilevel
population.

\hypertarget{finite-visibility-and-reference-examples}{%
\section{Finite visibility and reference
examples}\label{finite-visibility-and-reference-examples}}

Put \(\ell_{G,m}(q)=[x^m]L_G(x,q)\). Splitting weak compositions by
their multiplicity patterns gives (SNG-F14)

\[
 \ell_{G,1}=\sum_iq^{d_i},\qquad
 \ell_{G,2}=\sum_iq^{4d_i}+\sum_{i<j}q^{d_i+d_j-2A_{ij}},
\]

\[
 \begin{aligned}
 \ell_{G,3}={}&\sum_iq^{9d_i}
 +\sum_{i\ne j}q^{4d_i+d_j-4A_{ij}}\\
 &+\sum_{i<j<t}q^{d_i+d_j+d_t-2(A_{ij}+A_{it}+A_{jt})}.
 \end{aligned}
\]

These formulas also delimit the meaning of a small-depth witness. Volume
one observes degrees. With degrees fixed, volume two observes the edge
degree-sum polynomial \(M_G(q)=\sum_{ij\in E}q^{d_i+d_j}\), since then
\(\ell_{G,2}-\ell_{H,2}=(q^{-2}-1)(M_G-M_H)\). The current volume-three
witness separates an aggregated binary census by an ordered two-vertex
height pattern. It is not a claim that the graph energy contains a new
three-body interaction.

All reference inputs are explicit. On vertices \(0,\ldots,6\), the first
pair has edge sets

\[
 \begin{aligned}
 E(G_2)&=\{01,02,03,04,12,35,36\},\\
 E(H_2)&=\{01,02,03,04,15,16,56\}.
 \end{aligned}
\]

Their degree multisets and Laplacian characteristic polynomials agree.
The latter, computed exactly in integer arithmetic, is

\[
 t^7-14t^6+73t^5-178t^4+209t^3-112t^2+21t.
\]

The unique degree-four vertex lies in a triangle in the first graph and
not the second, so they are nonisomorphic (SNG-E001). Direct coefficient
enumeration gives the first difference

\[
 \ell_{G_2,2}-\ell_{H_2,2}=2q^2(1-q)^3(1+q)=2q^2(1-q)^2(1-q^2).
\]

Its source is elementary: the degree-sum polynomials differ,
\(M_{G_2}-M_{H_2}=2q^4(1-q)^2\). The pair therefore separates the height
population from the Laplacian spectrum through degree mixing, not
through a higher-order relation.

The second pair has edge sets

\[
 \begin{aligned}
 E(G_3)&=\{03,04,06,12,13,23,24,35\},\\
 E(H_3)&=\{01,04,06,15,16,26,36,45\}.
 \end{aligned}
\]

Its graphs have the same complete binary cut-volume polynomial:
enumeration of all 128 subsets of each graph agrees at every pair of
subset size and boundary size. Their unique degree-four vertices have
respectively one and two leaf neighbors, proving nonisomorphism
(SNG-E002). The local coefficients agree through volume two, while

\[
 \ell_{G_3,3}-\ell_{H_3,3}
 =q^3(1-q)(1-q^3)^2(1-q^4)^2.
\]

This entire discrepancy occurs in span two. Span one agrees by the
binary-census equality, and span three agrees by the equal degree
multisets. The degree-sum polynomials agree, which explains agreement at
volume two. The finer joint degree table of the edges differs, and the
unequal-height sector weights endpoint degrees through \(4d_i+d_j\).
Both differences are strictly positive on the open interval between zero
and one, as their factorizations show.

The elementary invariants shared by each pair are summarized below. The
data are in \texttt{results/exact\_examples.json}.

\begin{longtable}[]{@{}
  >{\raggedright\arraybackslash}p{(\columnwidth - 12\tabcolsep) * \real{0.1364}}
  >{\raggedleft\arraybackslash}p{(\columnwidth - 12\tabcolsep) * \real{0.1818}}
  >{\raggedright\arraybackslash}p{(\columnwidth - 12\tabcolsep) * \real{0.1364}}
  >{\raggedright\arraybackslash}p{(\columnwidth - 12\tabcolsep) * \real{0.1364}}
  >{\raggedright\arraybackslash}p{(\columnwidth - 12\tabcolsep) * \real{0.1364}}
  >{\raggedright\arraybackslash}p{(\columnwidth - 12\tabcolsep) * \real{0.1364}}
  >{\raggedright\arraybackslash}p{(\columnwidth - 12\tabcolsep) * \real{0.1364}}@{}}
\toprule\noalign{}
\begin{minipage}[b]{\linewidth}\raggedright
Pair
\end{minipage} & \begin{minipage}[b]{\linewidth}\raggedleft
Order
\end{minipage} & \begin{minipage}[b]{\linewidth}\raggedright
Degree multiset
\end{minipage} & \begin{minipage}[b]{\linewidth}\raggedright
Laplacian char. polynomial
\end{minipage} & \begin{minipage}[b]{\linewidth}\raggedright
Binary cut census
\end{minipage} & \begin{minipage}[b]{\linewidth}\raggedright
Edge degree-sum polynomial
\end{minipage} & \begin{minipage}[b]{\linewidth}\raggedright
First difference: volume, span
\end{minipage} \\
\midrule\noalign{}
\endhead
\bottomrule\noalign{}
\endlastfoot
\(G_2,H_2\) & 7 & equal & equal & different & different & 2, span 1 \\
\(G_3,H_3\) & 7 & equal & different & equal & equal & 3, span 2 \\
cube, Möbius ladder & 8 & equal (3-regular) & different & different &
equal (regular) & 4, span 1 \\
\end{longtable}

\textbf{Kernel-checked scope.} For the first two pairs, the decisive
coefficients are checked in Lean 4.12.0 (core library only, no project
axioms) by kernel evaluation of the released certificates
\texttt{results/certificates/spectral.json} and
\texttt{results/certificates/binary\_census.json}. Write X for
\texttt{spectral} or \texttt{binary} in namespace
\texttt{Singularity.Cert}. The following hold for the stated edge lists
on seven vertices.

\begin{itemize}
\tightlist
\item
  \texttt{X\_local} and \texttt{XCmp\_local}: every coefficient
  \([x^Vq^e]L\) of each graph, for all V at most 4 and all integers e,
  equals the transcribed table. The coefficient is defined semantically
  by \texttt{allCount}, as a count of all length-7 profiles of volume V
  and energy e.
\item
  \texttt{X\_normalized}: the minimum-zero coefficients \([x^Vq^e]B\) of
  the certificate graph \(G_2\), respectively \(G_3\), through volume 4.
  This follows from P07 because these volumes are below 7.
\item
  \texttt{X\_agree\_below}: the two graphs agree at every volume below
  the first difference.
\item
  \texttt{X\_first\_difference}, \texttt{X\_factorization} and
  \texttt{X\_discrepancy}: the first-difference polynomial is the
  displayed one, it equals the expansion of the factorization
  \(2q^2(1-q)^2(1-q^2)\), respectively \(q^3(1-q)(1-q^3)^2(1-q^4)^2\),
  and it is nonzero.
\end{itemize}

Rejection theorems show that one changed edge or one changed decisive
coefficient is refused. The trusted boundary is the transcription script
\texttt{formal/tools/gen\_certificates.py}, which copies the JSON
numbers into Lean literals and records the input SHA-256 digests. These
are \texttt{379d5c47…ec6b2} for \texttt{spectral.json} and
\texttt{e654012d…e4b2} for \texttt{binary\_census.json}; full digests
are in \texttt{formal/README.md}. The script's \texttt{-\/-check} mode
detects a stale transcription.

The kernel does not check the span-marked rows, which include the
span-two localization above. It also does not check the Euler exponents,
the binary cut census, the Laplacian polynomials, nonisomorphism,
simplicity of the edge lists, or anything about the Python
implementation. Those rest on exact software replay, on two independent
enumerators, and on the short arguments given here. The package builds
from a clean download of the pinned Lean release in the project's
continuous integration, which also audits that the listed theorems
depend only on Lean's standard axioms; the build command and theorem
list are in \texttt{formal/README.md}.

A third pair, the three-dimensional cube and the eight-vertex Möbius
ladder, consists of triangle-free cubic graphs. They are nonisomorphic
because the cube is bipartite and the ladder contains a five-cycle.
Their local coefficients agree through volume three and first differ at
volume four by \(2q^4(1-q^2)^4\). That difference lies in span one, the
binary sector. The example illustrates delayed visibility after matching
regular low-order data. It does not duplicate the binary-census equality
of the second pair. Its coefficients are confirmed by the independent
software replay of
\texttt{results/certificates/regular\_cube\_mobius.json}. They are not
kernel-checked.

\hypertarget{arithmetic-assembly-and-the-comparison-theorem}{%
\section{Arithmetic assembly and the comparison
theorem}\label{arithmetic-assembly-and-the-comparison-theorem}}

For a positive integer \(n\), let \(v_p(n)\) be its \(p\)-adic
valuation. Define

\[
 \mathcal E_G(n_1,\ldots,n_k)
 =\sum_pE_G(v_p(n_1),\ldots,v_p(n_k)),
\]

\[
 Z_G(s,q)=\sum_{n_1,\ldots,n_k\ge1}
 \frac{q^{\mathcal E_G(\mathbf n)}}{(n_1\cdots n_k)^s}.
\]

For fixed q in the closed unit interval and real part of s greater than
one, absolute convergence and unique factorization give
\(Z_G(s,q)=\prod_pL_G(p^{-s},q)\) (SNG-A01). The construction is
multiplicative across primes, while each prime's local factor retains
the graph interaction among its exponent coordinates.

\textbf{Theorem 3 (first-discrepancy visibility; SNG-A11).} Fix q in the
closed interval {[}0,1{]}. Suppose two graphs with the same number of
vertices satisfy \(\ell_{G,m}(q)=\ell_{H,m}(q)\) for m below r, and
\(\delta=\ell_{G,r}(q)-\ell_{H,r}(q)\ne0\). Then their ratio Euler
product converges for real s greater than 1/r and equals

\[
 R_{G,H}(s,q)=\zeta(rs)^\delta K_{G,H}(s,q),
\]

where K has a finite positive limit at s=1/r and is holomorphic and
nonzero in a complex neighborhood of that point. Consequently

\[
 R_{G,H}(s,q)\sim K_{G,H}(1/r,q)(rs-1)^{-\delta}
 \quad(s\downarrow1/r).
\]

\textbf{Proof.} Constant terms are one, so cancellation of the lower
local coefficients gives \(L_G/L_H=1+\delta x^r+O(x^{r+1})\). The ratio
Euler product therefore has an absolutely summable logarithm for real s
greater than 1/r. Every local factor is positive there. Multiplication
by \((1-x^r)^\delta\) removes the first discrepancy, yielding
\(1+O(x^{r+1})\). Thus

\[
 K_{G,H}(s,q)=\prod_p(1-p^{-rs})^\delta
              \frac{L_G(p^{-s},q)}{L_H(p^{-s},q)}
\]

converges for real s greater than 1/(r+1). Near s=1/r, finitely many
small-prime factors are nonzero by positivity and continuity; the
remaining tail has a normally convergent logarithm. This proves the
local complex statement. The zeta Euler product proves the factorization
on the right of 1/r, and \(\zeta(1+t)\sim1/t\) gives the asymptotic.
\(\square\)

The hypothesis concerns values at the chosen q. A polynomial discrepancy
can vanish at an endpoint or at an interior root, which changes the
first visible volume at that parameter. For the three reference pairs
the polynomial differences are positive on 0\textless q\textless1, so
the theorem applies there with r=2, 3 and 4. At q=1/2 the exponents are
δ=3/32 at s=1/2, δ=11025/262144 at s=1/3, and δ=81/2048 at s=1/4 for the
cube comparison. These are exact rational evaluations of the displayed
differences. Nonintegral exponents give power or branch singularities,
not ordinary poles. The quotient is initially a quotient of convergent
positive series on s\textgreater1; its ratio Euler product supplies the
extension used in the theorem. No global meromorphic-continuation claim
follows. Theorem 3 is a human-readable proof; it is not formalized.

\textbf{Certified regular factors.} The constant \(K_{G,H}(1/r,1/2)\)
has been enclosed rigorously for the two seven-vertex pairs (SNG-F38,
\texttt{results/certified\_k.json}). The enclosure assumes that both
graphs are connected, q is rational in (0,1) and s=a/b is rational with
(r+1)s\textgreater1. It keeps all primes up to a cutoff \(P=N^b\)
exactly and bounds the omitted primes by an explicit constant times
\(P^{1-(r+1)s}/((r+1)s-1)\). Truncated heights are bounded by a
quadratic tail estimate (SNG-F39). The computation uses 256-bit
fixed-point arithmetic with outward rounding and series enclosures of
the logarithms. The retained intervals, rounded outward from their exact
dyadic endpoints, are as follows.

\begin{longtable}[]{@{}
  >{\raggedright\arraybackslash}p{(\columnwidth - 6\tabcolsep) * \real{0.2308}}
  >{\raggedright\arraybackslash}p{(\columnwidth - 6\tabcolsep) * \real{0.2308}}
  >{\raggedleft\arraybackslash}p{(\columnwidth - 6\tabcolsep) * \real{0.3077}}
  >{\raggedright\arraybackslash}p{(\columnwidth - 6\tabcolsep) * \real{0.2308}}@{}}
\toprule\noalign{}
\begin{minipage}[b]{\linewidth}\raggedright
Pair
\end{minipage} & \begin{minipage}[b]{\linewidth}\raggedright
\(s\)
\end{minipage} & \begin{minipage}[b]{\linewidth}\raggedleft
Primes kept
\end{minipage} & \begin{minipage}[b]{\linewidth}\raggedright
Enclosure of \(K(s,1/2)\)
\end{minipage} \\
\midrule\noalign{}
\endhead
\bottomrule\noalign{}
\endlastfoot
\(G_2,H_2\), \(r=2\), \(\delta=3/32\) & \(1/2\) & 1229 (\(p\le10^4\)) &
\([0.84973,\,0.85750]\) \\
\(G_2,H_2\) & \(3/5\) & 985 (\(p\le6^5\)) & \([0.89931,\,0.89971]\) \\
\(G_2,H_2\) & \(1\) & 168 (\(p\le10^3\)) &
\([0.9678188,\,0.9678191]\) \\
\(G_3,H_3\), \(r=3\), \(\delta=11025/262144\) & \(1/3\) & 2961
(\(p\le30^3\)) & \([0.90084,\,0.94230]\) \\
\end{longtable}

Thus \(R_{G_2,H_2}(s,1/2)\sim C_2(2s-1)^{-3/32}\) with
\(C_2\in[0.84973,0.85750]\). Likewise
\(R_{G_3,H_3}(s,1/2)\sim C_3(3s-1)^{-11025/262144}\) with
\(C_3\in[0.90084,0.94230]\). The exponents come from Theorem 3 and the
exact coefficients; they are not fitted. The binary interval is wide
because the omitted-prime bound decays only like \(P^{-1/3}\) at s=1/3.
These are computer-assisted interval computations in Python. They are
not kernel-checked. No enclosure has been computed for the cube
comparison.

\hypertarget{synchronization-and-its-singular-limit}{%
\section{Synchronization and its singular
limit}\label{synchronization-and-its-singular-limit}}

Connectedness implies zero energy exactly at constant heights. Hence
\(Z_G(s,0)=\zeta(ks)\), whereas \(Z_G(s,1)=\zeta(s)^k\) (SNG-A04). More
generally, the first local coefficient is \(\alpha_G(q)=\sum_iq^{d_i}\).
For fixed q in {[}0,1{]}, removal of the linear Euler term gives
(SNG-A09)

\[
 Z_G(s,q)=\zeta(s)^{\alpha_G(q)}H_G(s,q),
\]

initially for real part of s above one, with H holomorphic for real part
of s above 1/2 and positive and nonzero near real s=1. Every fixed
positive q therefore restores an s=1 divergence, while exact
synchronization q=0 leaves the finite value \(\zeta(k)\) there. The
order of these two limits matters.

Let d be minimum degree and b the number of vertices attaining it.
Suppose q tends to zero and real s tends down to one, along sequences,
with \(q^d\log(1/(s-1))\to c<\infty\). Then (SNG-A16)

\[
 \frac{Z_G(s,q)}{\zeta(ks)}\longrightarrow e^{bc}.
\]

To justify this limit, use translation-normalized local factors and
subtract their common linear term. For real s at least one, each such
logarithm is bounded by an explicit constant times
\(q^\lambda p^{-2s}\), where λ is the edge connectivity. The bound comes
from the weak-composition coefficient count and the fact that every
nonzero minimum-zero profile has energy at least λ (SNG-A15). Summing
over primes, the normalized product tends to one as q tends to zero,
uniformly in s. Finally \(\alpha_G(q)=bq^d+O(q^{d+1})\) and
\(\log\zeta(s)=\log(1/(s-1))+O(s-1)\) supply the exponential limit. This
argument describes a boundary layer in the declared real domain,
independently of any global continuation question.

The first positive energy has a separate combinatorial interpretation.
Integer coarea bounds energy below by the sum of threshold cut sizes,
hence by λ times span. Equality at the minimum positive energy forces a
binary minimum-cut indicator (SNG-A05). Consequently, initially for real
part of s above one, the coefficient of \(q^\lambda\) in
\(\log(Z_G/\zeta(ks))\) is the sum of prime-zeta functions
\(\sum_{|\partial S|=\lambda}P(|S|s)\), counting both cut orientations
(SNG-A06). This links cut side sizes to an analytic coefficient without
claiming that minimum cuts identify the whole graph.

\hypertarget{computational-and-formal-scope}{%
\section{Computational and formal
scope}\label{computational-and-formal-scope}}

The released computational core uses integer and rational arithmetic and
the Python standard library. It supplies weak-composition counts, marked
sectors, exact Laplacian characteristic polynomials, formal Euler
exponents with independent recomposition, constrained original-label
rank and unrank, forest message passing, root-of-unity residue
numerators, weighted hard contraction and primitive-profile inversion. A
finite certificate records a graph source, coefficient range and
declared comparison. Replay parses it strictly and recomputes all
coefficients by two routes that do not reuse the builder's enumerator:
stars-and-bars compositions and level chains. It also checks the
normalization and threshold identities.
\texttt{PYTHONPATH=src\ python\ scripts/reproduce.py\ -\/-check} replays
the four retained certificates and the three comparison pairs. These are
software-level exact checks relative to the Python interpreter.

The Lean package in \texttt{formal/} crosses the remaining boundary only
where stated. It kernel-checks the coefficientwise normalization
identity of Proposition 1, the nested energy identity of Proposition 2,
and the decisive unmarked coefficient tables of the first two pairs, as
listed in Section 3. It also kernel-checks the finite cancellation
behind Theorem 3. A first coefficient discrepancy at volume r gives the
quotient \(1+\delta x^r\) modulo \(x^{r+1}\) for any series over any
commutative ring, and the binomial factor cancels it. For the first two
pairs the quotients are \(1+2q^2(1-q)^2(1-q^2)x^2\) and
\(1+q^3(1-q)(1-q^3)^2(1-q^4)^2x^3\). The analytic conclusion appears in
Lean only as a conditional theorem. Its real-analysis premises (the
Euler product, the pole of zeta, convergence of products and the local
series expansions) are stated, not proved. The Lean package does not
formalize the span-marked series, the prime products over the real
numbers or any unconditional analytic statement of Sections 4 and 5. It
is not linked to the Python implementation.

An exhaustive enumeration covers all labeled simple graphs on one to six
vertices, 26,704 of which are connected on six vertices. Connected
graphs were grouped into isomorphism classes by an exact minimal
adjacency code over degree-preserving relabelings. This yields
respectively 1, 1, 2, 6, 21 and 112 classes
(\texttt{results/small\_census.json},
\texttt{scripts/census.py\ -\/-max-order\ 6\ -\/-max-volume\ 6}). For
order six, exact polynomial local coefficients through volumes one, two
and three distinguish respectively 68, 100 and 112 classes. In this
population the binary cut census also separates all classes, so a
connected nonisomorphic pair with equal binary census, such as the
second pair, needs at least seven vertices. This is a complete finite
census with an explicit generation boundary. It does not imply
completeness at larger orders. In particular, the seven- and
eight-vertex examples illustrate why visibility depth must remain part
of the specification.

\hypertarget{open-structural-questions}{%
\section{Open structural questions}\label{open-structural-questions}}

The main inverse problem asks whether the fully marked series identifies
a graph in a declared class. A negative answer should describe a
transformation or an exact pair preserving all heights, rather than
merely a long agreeing prefix. A complementary question asks whether
nonisomorphic pairs can have first-separation depth tending to infinity.
Both questions concern scalar marks. If every vertex receives an
independent variable, labeled adjacency is recovered already from the
one- and two-coordinate coefficients (SNG-F18), making that unrestricted
version much easier.

The prime lift adds an analytic hierarchy to these local questions. It
does not add graph information: coefficients at prime powers recover the
local coefficients, and local coefficients determine the Euler product.
Its contribution is to turn a finite first discrepancy into a controlled
singularity, allowing relational depth and arithmetic location to be
studied together.

\hypertarget{references}{%
\section*{References}\label{references}}
\addcontentsline{toc}{section}{References}

{[}BHN1997{]} R. Bacher, P. de la Harpe and T. Nagnibeda, \emph{The
lattice of integral flows and the lattice of integral cuts on a finite
graph}, Bulletin de la Societe Mathematique de France 125 (1997),
167-198. DOI: 10.24033/bsmf.2303.

{[}AM2009{]} D. Andren and K. Markstrom, \emph{The bivariate Ising
polynomial of a graph}, Discrete Applied Mathematics 157 (2009),
2515-2524. DOI: 10.1016/j.dam.2009.02.021.

{[}BPR2024{]} R. Bauerschmidt, J. Park and P.-F. Rodriguez, \emph{The
Discrete Gaussian model, I. Renormalisation group flow at high
temperature}, Annals of Probability 52 (2024), no. 4, 1253-1359. DOI:
10.1214/23-AOP1658. arXiv:2202.02286.

{[}Moree2005{]} P. Moree, \emph{The formal series Witt transform},
Discrete Mathematics 295 (2005), 143-160. DOI:
10.1016/j.disc.2005.03.004.

\end{document}
